\textbf{Colors and the zero contour.} Red means positive mean $\rho_e$: the degree-10 model has a larger clean-target error than the degree-2 model despite having no larger optimal training error. Blue means negative mean $\rho_e$, so degree 10 generalizes better in this particular comparison. White is near zero and the black contour marks the empirical zero-mean boundary. A negative value does not prove the degree-10 model cannot overfit relative to another model.

Colors saturate at $\pm0.25$ for visibility. The actual complexity-grid means range from $-0.296410$ to $1.032363\times10^6$, and noise-grid means from $-0.146917$ to $4.812316\times10^6$. These values were not clipped. Some finite-sample Legendre designs are poorly conditioned: the smallest recorded degree-10 design singular value is $1.304\times10^{-5}$. Large fitted coefficients can therefore amplify noise or model misspecification into large population errors, particularly where sample coverage is poor. All trials, including these extreme ones, contribute to the required arithmetic mean.

\textbf{Sample count.} More observations generally reduce the estimation instability of the flexible model in the observed grids and allow its approximation advantage to emerge. For example, at $Q_g=20$, $\sigma^2=0.1$, the mean error difference changes from about $968028.958439$ at $n=20$, to $11.581522$ at $n=40$, to $2.483317$ at $n=60$, and to $-0.119340$ at $n=120$. At $n=120$, the mean risks are $0.409165$ for degree 2 and $0.289825$ for degree 10. This is an observed tendency, not a theorem of monotonicity for each sample or every grid cell.

\textbf{Noise variance.} Noise can be fitted by the extra coefficients of $f_{10}$ and increase estimation variance. At $Q_g=20,n=120$, the means below move from a degree-10 advantage to a disadvantage as noise increases:
\begin{center}\small
\begin{tabular}{@{}rrrrr@{}}\toprule
$\sigma^2$ & $0$ & $0.1$ & $1$ & $2$\\\midrule
Mean $\rho_e$ & $-0.146917$ & $-0.119340$ & $0.037381$ & $0.196897$\\\bottomrule
\end{tabular}\end{center}
Finite paired Monte Carlo estimates need not be monotone at every cell. For example, at $n=40$ the zero-noise mean is $11.665551$ and the variance-$0.1$ mean is $11.581522$. Sample-specific cross terms can make such local comparisons nonmonotone; they are not grounds for discarding trials or changing the prescribed mean.

\textbf{Target complexity.} Increasing $Q_g$ is not always harmful. At $\sigma^2=0.1,n=120$, the mean $\rho_e$ is $0.014199$ for $Q_g=1$, $-0.071197$ for $Q_g=3$, $-0.296410$ for $Q_g=10$, and $-0.119340$ for $Q_g=20$. For targets of degree at most 2, both model classes contain the target and the extra coefficients mainly add noise-fitting variance. For intermediate degrees up to 10, the larger model can represent structure that the quadratic omits. Above degree 10, both models are misspecified; estimation error, omitted high-order components and the degree-dependent target normalization all matter. Consequently target degree alone does not determine the relative error.

\textbf{Can positive relative overfitting occur without observation noise?} Yes. The noise-grid mean at $\sigma^2=0,Q_g=20,n=20$ is $558740.582041>0$, whereas at $n=120$ it is $-0.146917$. With degree 20 the target is outside both fitted model classes, and fitting its values at a limited random set of inputs can induce a poor degree-10 approximation elsewhere, especially when the design is ill-conditioned. Observation noise is therefore not necessary. If the noiseless target were contained in the degree-10 class and the design had full column rank, an exact fit would instead recover it in exact arithmetic; this positive-error argument concerns the misspecified target case.

\textbf{Limits of the contour.} The displayed boundary is an empirical comparison of two models under this fixed random bank. The grids share random numbers, so adjacent cells are paired rather than independent. Within each cell there are 1000 independent trials, but rare large errors can strongly affect the mean: at $Q_g=20,\sigma^2=0.1,n=20$, its standard error is about $498426.190234$. The black line is not evidence of a universal sharp transition, nor a guarantee about an individual fitted model.
